\section*{Capítulo \ref{ch:g12:work-and-energy} --- Trabajo y energía mecánica}

\begin{solution}{exo:g12:work-and-energy:1}
$E_p = \tfrac12 \times 200 \times 0.050^2 = \qty{0.25}{J}$. Con la
elongación duplicada: $\times 4$, es decir, \qty{1.0}{J}. Los segundos
\qty{5.0}{cm}, por sí solos, cuestan $1.0 - 0.25 = \qty{0.75}{J}$ ---
allí la \omterm{def:g10:inertia:force}{fuerza} es mayor.
\end{solution}

\begin{solution}{exo:g12:work-and-energy:2}
$W = mg(z_A - z_B) = 0.90 \times 9.81 \times (12 - 30) \approx
\qty{-1.6e2}{J}$, sea cual sea el recorrido --- lo mismo por la subida
recta. Ida y vuelta: \qty{0}{J} (\omterm{def:g12:work-and-energy:conservative}{fuerza conservativa}).
\end{solution}

\begin{solution}{exo:g12:work-and-energy:3}
(a) y (b), \omterm{def:g12:work-and-energy:conservative}{conservativas}: el \omterm{def:g11:work-of-force:work}{trabajo} queda fijado por los extremos
(desnivel; elongación). (c) y (e), \omterm{def:g12:work-and-energy:conservative}{no conservativas}: $-fL$ y el
\omterm{def:g10:inertia:inventory}{rozamiento} del aire cobran por \omterm{def:g10:orders-of-magnitude:unit}{metro} recorrido. (d) no trabaja: es
siempre perpendicular al movimiento.
\end{solution}

\begin{solution}{exo:g12:work-and-energy:4}
$h = L(1 - \cos\qty{60}{\degree}) = \qty{1.0}{m}$;
$v = \sqrt{2 \times 9.81 \times 1.0} \approx \qty{4.4}{m/s}$. La tensión
es perpendicular al movimiento en todo instante: \omterm{def:g11:work-of-force:power}{potencia} nula, \omterm{def:g11:work-of-force:work}{trabajo}
nulo.
\end{solution}

\begin{solution}{exo:g12:work-and-energy:5}
$F = P/v = 250/9.0 \approx \qty{28}{N}$. A \qty{5.0}{m/s}:
$F = 250/5.0 = \qty{50}{N}$ (sobre todo la componente de la gravedad a
lo largo de la pendiente).
\end{solution}

\begin{solution}{exo:g12:work-and-energy:6}
$E_p = \tfrac12 \times 450 \times 0.080^2 = \qty{1.44}{J}$;
$v = \sqrt{2 \times 1.44/0.020} = \qty{12}{m/s}$;
$h = E_p/mg = 1.44/(0.020 \times 9.81) \approx \qty{7.3}{m}$.
\end{solution}

\begin{solution}{exo:g12:work-and-energy:7}
Trapecios (pasos de \qty{2.0}{cm}):
$(0.03 + 0.10 + 0.19 + 0.30) = \qty{0.62}{J}$. Modelo de recta:
$k = 18.0/0.080 = \qty{225}{N/m}$, $\tfrac12 kx^2 = \qty{0.72}{J}$. La
curva medida se hunde por debajo de la recta para $x$ pequeña, así que
el área verdadera es menor.
\end{solution}

\begin{solution}{exo:g12:work-and-energy:8}
$\Delta E_m = \tfrac12 \times 55 \times 9.0^2 - 55 \times 9.81 \times
8.0 = 2228 - 4316 \approx \qty{-2.1e3}{J}$;
$f = 2089/25 \approx \qty{84}{N}$. El \omterm{def:g10:universal-gravitation:weight}{peso} es conservativo y la \omterm{def:g12:newtons-laws:friction}{reacción
normal} no trabaja: solo entran el desnivel y la \emph{longitud} del
camino, no su forma.
\end{solution}

\begin{solution}{exo:g12:work-and-energy:9}
$H_{\min} = \tfrac52 R = \qty{0.50}{m}$. Desde \qty{60}{cm}:
$v_{\text{cima}} = \sqrt{2 \times 9.81 \times (0.60 - 0.40)} \approx
\qty{2.0}{m/s}$, y $v_{\text{cima}}^2/(gR) = 2.0$ --- el doble del
mínimo. El \omterm{def:g10:inertia:inventory}{rozamiento} va rascando \omterm{def:g10:energy-conservation:energy}{energía} todo el camino, así que las
pistas reales necesitan altura de más.
\end{solution}

\begin{solution}{exo:g12:work-and-energy:10}
\omterm{def:g12:work-and-energy:diagram}{Puntos de retorno}: $8.0\,x^2 = 2.0$, $x = \pm\qty{0.50}{m}$. Rapidez
máxima en $x = 0$: $v = \sqrt{2 \times 2.0/0.100} \approx
\qty{6.3}{m/s}$. Una oscilación de ida y vuelta en un pozo parabólico
--- la $E_p$ del \omterm{prop:g12:oscillators-and-time:spring-period}{oscilador masa--resorte} tiene exactamente esta forma.
\end{solution}

\begin{solution}{exo:g12:work-and-energy:11}
$\tfrac12 kx^2 = \tfrac12 mv^2$:
$x = v\sqrt{m/k} = 2.0\sqrt{1200/\num{6.0e5}} \approx \qty{8.9}{cm}$. A
\qty{4.0}{m/s}: \qty{18}{cm} --- el doble, no el cuádruple: la \omterm{def:g10:energy-conservation:energy}{energía}
se cuadruplica, pero $E_p$ crece como $x^2$, así que $x$ solo se
duplica.
\end{solution}

\begin{solution}{exo:g12:work-and-energy:12}
(a) $E_m = \qty{5.0}{J} < \qty{6.0}{J}$: no puede pasar; se para en la
ladera donde $E_p = \qty{5.0}{J}$ y vuelve rodando. (b) En la cima:
$E_k = \qty{1.0}{J}$, $v = \sqrt{2 \times 1.0/0.50} = \qty{2.0}{m/s}$;
en el valle: $E_k = \qty{5.0}{J}$, $v \approx \qty{4.5}{m/s}$.
(c) $x = 0$: inestable (máximo); $x = \qty{1.5}{m}$: estable (mínimo).
\end{solution}

\begin{solution}{exo:g12:work-and-energy:13}
$\tfrac12 mv^2 = mgR_E$:
$v = \sqrt{2 \times 9.81 \times \num{6.37e6}} \approx \qty{1.12e4}{m/s}
= \qty{11.2}{km/s}$. Todos los términos son proporcionales a $m$: da
igual sonda que guijarro. Justo por debajo, el lanzamiento sube aún
muchísimo, encuentra un \omterm{def:g12:work-and-energy:diagram}{punto de retorno} y vuelve a caer: queda ligado.
\end{solution}

\begin{solution}{exo:g12:work-and-energy:14}
$E_m = \tfrac12 kA^2 = \tfrac12 \times 40 \times 0.060^2 =
\qty{7.2e-2}{J}$;
$v_{\max} = \sqrt{2E_m/m} = \sqrt{2 \times 0.072/0.250} \approx
\qty{0.76}{m/s}$; $E_k = E_p$ donde $\tfrac12 kx^2 = E_m/2$:
$x = \pm A/\sqrt 2 \approx \pm\qty{4.2}{cm}$.
\end{solution}

\begin{solution}{exo:g12:work-and-energy:15}
$h = v^2/2g = 9.5^2/19.62 \approx \qty{4.6}{m}$; listón
$\approx 4.6 + 1.0 = \qty{5.6}{m}$. El récord es mayor: los saltadores
añaden \omterm{def:g11:work-of-force:work}{trabajo} con brazos y piernas sobre la pértiga flexionada. Pero la
estimación energética fija la escala --- ningún salto de velocista
llegará a \qty{10}{m}.
\end{solution}

\begin{solution}{pb:g12:work-and-energy:1}
\textbf{1.} $v = \sqrt{2gH} = \sqrt{2 \times 9.81 \times 45} \approx
\qty{29.7}{m/s}$: \omterm{def:g10:universal-gravitation:weight}{peso} conservativo, \omterm{def:g12:newtons-laws:friction}{reacción normal} sin \omterm{def:g11:work-of-force:work}{trabajo} --- el
perfil desaparece.

\textbf{2.} Puerta: $E_m = mgH = 70 \times 9.81 \times 45 \approx
\qty{3.09e4}{J}$. Despegue: $E_m = \tfrac12 \times 70 \times 26^2
\approx \qty{2.37e4}{J}$.

\textbf{3.} $\Delta E_m \approx \qty{-7.2e3}{J}$; el \omterm{thm:g12:work-and-energy:emtheorem}{teorema de la
energía mecánica} se lo carga a las \omterm{def:g12:work-and-energy:conservative}{fuerzas no conservativas}: el
\omterm{def:g10:inertia:inventory}{rozamiento} de la nieve y el del aire.

\textbf{4.} $f = \num{7.2e3}/90 \approx \qty{80}{N}$.

\textbf{5.} $\num{7.2e3}/\num{3.09e4} \approx 23\,\%$. La cera reduce
$f$ y agacharse reduce el \omterm{def:g10:inertia:inventory}{rozamiento} del aire: ambas cosas encogen $fL$,
el único término negociable.

\textbf{6.} \omterm{def:g12:projectile-motion:freefall}{Caída libre} con velocidad inicial horizontal --- un tiro
parabólico. $v_x = \qty{26}{m/s}$;
$v_y = \sqrt{2 \times 9.81 \times 40} \approx \qty{28.0}{m/s}$.

\textbf{7.} $v = \sqrt{26^2 + 28.0^2} = \sqrt{1461} \approx
\qty{38.2}{m/s}$; por \omterm{def:g10:energy-conservation:energy}{energía}, $v = \sqrt{v_0^2 + 2gh} =
\sqrt{676 + 785}$ --- idéntico. Unos \qty{138}{km/h}.

\textbf{8.} $t = v_y/g = 28.0/9.81 \approx \qty{2.9}{s}$;
$d = v_0 t \approx \qty{74}{m}$.

\textbf{9.} No. La sustentación es perpendicular a la velocidad (no
trabaja) y el \omterm{def:g10:inertia:inventory}{rozamiento} se le opone (\omterm{def:g11:work-of-force:work}{trabajo} negativo):
$\Delta E_m \leq 0$, así que la rapidez real de recepción solo puede ser
\emph{menor} --- el aire alarga el vuelo, no lo alimenta.

\textbf{10.} $\alpha = \arctan(28.0/26) \approx \qty{47}{\degree}$ por
debajo de la horizontal.

\textbf{11.} Ángulo con la pendiente: $47 - 35 = \qty{12}{\degree}$.
$v_\parallel = 38.2\cos\qty{12}{\degree} \approx \qty{37.4}{m/s}$;
$v_\perp = 38.2\sin\qty{12}{\degree} \approx \qty{7.9}{m/s}$.

\textbf{12.} Absorbida: $\tfrac12 \times 70 \times 7.9^2 \approx
\qty{2.2e3}{J}$, frente a $\tfrac12 \times 70 \times 38.2^2 \approx
\qty{5.1e4}{J}$ en total: alrededor del $4\,\%$.

\textbf{13.} $h_{\text{eq}} = 7.9^2/19.62 \approx \qty{3.2}{m}$ --- un
salto atrevido desde un muro. En llano:
$38.2^2/19.62 \approx \qty{74}{m}$ --- no se sobrevive.

\textbf{14.} Una pendiente paralela a la \omterm{def:g10:relative-motion:trajectory}{trayectoria} de vuelo mantiene
$v_\perp$ pequeña, así que las piernas absorben unos \omterm{def:g10:orders-of-magnitude:unit}{metros} y no la
montaña entera de \omterm{def:g11:mechanical-energy:kinetic}{energía cinética}.

\textbf{15.} $H'_0 = v_0'^2/2g = 576/19.62 \approx \qty{29.4}{m}$.

\textbf{16.} $mgH' - f\,\dfrac{H'}{\sin\qty{35}{\degree}} =
\tfrac12 m v_0'^2$.

\textbf{17.} $H' = \dfrac{\tfrac12 \times 70 \times 576}
{686.7 - 80/0.574} = \dfrac{\num{20160}}{547} \approx \qty{37}{m}$.

\textbf{18.} $36.8/29.4 \approx 1.25$: un recargo del $25\,\%$.

\textbf{19.} Para \qty{60}{kg}:
$H' = \num{17280}/(588.6 - 139.5) \approx \qty{38.5}{m}$. Es la
saltadora \emph{más ligera} la que necesita la torre más alta: la
factura de \omterm{def:g10:inertia:inventory}{rozamiento} $fL$ es la misma, pero su presupuesto energético
$mgH'$ es menor.

\textbf{20.} Anuncia una torre de \qty{37}{m} para \qty{24}{m/s} de
despegue --- y confiesa que una cuarta parte de ella, unos \qty{7}{m} de
hormigón, es el recargo que el \omterm{def:g10:inertia:inventory}{rozamiento} le cobra al sueño sin
\omterm{def:g10:inertia:inventory}{rozamiento}.
\end{solution}
